\documentclass{article}
\usepackage{amsmath, amssymb}

\begin{document}

\section*{Pushing--Medium Gravity: Stress--Energy Tensor v1}

This document presents the stress--energy tensor of the Pushing--Medium (PM)
medium derived from the PM-fluid Lagrangian. The tensor describes the energy,
momentum, and stresses carried by the scalar field $\phi$ and the flow field
$\mathbf u$.

\section*{1. PM Lagrangian Density}

The PM-fluid Lagrangian on a flat background $\eta_{\mu\nu} =
\mathrm{diag}(-1,1,1,1)$ is


\[
\mathcal{L}
=
\frac{1}{2c_\phi^2}(\partial_t \phi)^2
- \frac{1}{2}|\nabla \phi|^2
+
\frac{1}{2c_u^2}|\partial_t \mathbf u|^2
- \frac{1}{2}|\nabla \mathbf u|^2
+
\phi\,\nabla\cdot(\phi\mathbf u)
+
\phi\,S_\phi
+
\mathbf u\cdot\mathbf S_u.
\]



For the stress--energy tensor we omit the explicit source terms
$\phi S_\phi$ and $\mathbf u\cdot\mathbf S_u$, which represent external
couplings.

\section*{2. Canonical Stress--Energy Tensor}

For fields $\Psi \in \{\phi, u_i\}$, the canonical tensor is


\[
T^{\mu}{}_{\nu}
=
\sum_{\Psi}
\frac{\partial \mathcal{L}}{\partial(\partial_\mu \Psi)}\,
\partial_\nu \Psi
- \delta^{\mu}{}_{\nu}\,\mathcal{L}.
\]



We use $x^0 = t$ and spatial indices $i,j = 1,2,3$.

\section*{3. Conjugate Momenta}

\subsection*{3.1 Scalar Field}



\[
\pi_\phi
=
\frac{\partial \mathcal{L}}{\partial(\partial_t \phi)}
=
\frac{1}{c_\phi^2}\,\partial_t \phi.
\]



\subsection*{3.2 Vector Field}



\[
\pi_{u_i}
=
\frac{\partial \mathcal{L}}{\partial(\partial_t u_i)}
=
\frac{1}{c_u^2}\,\partial_t u_i.
\]



\section*{4. Energy Density}

The energy density is


\[
\mathcal{E}
=
T^{00}
=
\pi_\phi\,\partial_t \phi
+
\pi_{u_i}\,\partial_t u_i
- \mathcal{L}.
\]



Substituting the momenta and Lagrangian (without sources):


\[
\mathcal{E}
=
\frac{1}{2c_\phi^2}(\partial_t \phi)^2
+ \frac{1}{2}|\nabla \phi|^2
+
\frac{1}{2c_u^2}|\partial_t \mathbf u|^2
+ \frac{1}{2}|\nabla \mathbf u|^2
- \phi\,\nabla\cdot(\phi\mathbf u).
\]



This is the total energy stored in the PM medium.

\section*{5. Momentum Density}

The momentum density is


\[
T^{0i}
=
\pi_\phi\,\partial_i \phi
+
\pi_{u_j}\,\partial_i u_j.
\]



Explicitly:


\[
T^{0i}
=
\frac{1}{c_\phi^2}(\partial_t \phi)(\partial_i \phi)
+
\frac{1}{c_u^2}(\partial_t u_j)(\partial_i u_j).
\]



\section*{6. Spatial Stress Tensor}

The spatial stress tensor is


\[
T^{ij}
=
\frac{\partial \mathcal{L}}{\partial(\partial_i \phi)}\,\partial^j \phi
+
\frac{\partial \mathcal{L}}{\partial(\partial_i u_k)}\,\partial^j u_k
- \delta^{ij}\,\mathcal{L}.
\]



From the gradient terms:


\[
\frac{\partial \mathcal{L}}{\partial(\partial_i \phi)}
= -\partial_i \phi,
\qquad
\frac{\partial \mathcal{L}}{\partial(\partial_i u_k)}
= -\partial_i u_k.
\]



Thus:


\[
T^{ij}
=
- (\partial_i \phi)(\partial^j \phi)
- (\partial_i u_k)(\partial^j u_k)
- \delta^{ij}\,\mathcal{L}_{\text{PM}},
\]


where


\[
\mathcal{L}_{\text{PM}}
=
\frac{1}{2c_\phi^2}(\partial_t \phi)^2
- \frac{1}{2}|\nabla \phi|^2
+
\frac{1}{2c_u^2}|\partial_t \mathbf u|^2
- \frac{1}{2}|\nabla \mathbf u|^2
+
\phi\,\nabla\cdot(\phi\mathbf u).
\]



\section*{7. Conservation Law}

In the absence of explicit sources:


\[
\partial_\mu T^{\mu\nu} = 0.
\]



This expresses local conservation of energy and momentum in the PM medium.

\section*{8. Physical Interpretation}

\begin{itemize}
    \item The terms $(\partial_t \phi)^2$ and $(\partial_t \mathbf u)^2$ represent kinetic energy of scalar and vector PM waves.
    \item The terms $|\nabla \phi|^2$ and $|\nabla \mathbf u|^2$ represent stored gradient energy, analogous to tension or pressure.
    \item The term $-\phi\,\nabla\cdot(\phi\mathbf u)$ represents compressive work done by the flow field on the scalar field.
    \item The momentum density $T^{0i}$ describes transport of PM energy by waves and flows.
    \item The stress tensor $T^{ij}$ describes directional tensions and pressures within the PM medium.
\end{itemize}

\section*{9. Summary}

The PM stress--energy tensor provides a complete description of energy,
momentum, and stresses in the PM medium. It is derived directly from the
PM-fluid Lagrangian and is consistent with the dynamic PM equations and the
effective metric structure.

\end{document}

